\begin{trapbox}
	\begin{description}[style=nextline, leftmargin=2.8em, labelsep=0.5em, itemsep=0.35em]
		\item[\textbf{\textcolor{CTrap}{易错点一（系数可以为零）}}] 当 $a=0$ 或 $b=0$ 时，认为不能用辅助角公式。
		\item[\textbf{\textcolor{CTrap}{正确理解（公式仍适用）：}}] 公式对任意不同时为零的 $a,b$ 都成立。若 $a=0$，则 $\cos\varphi=0$，$\sin\varphi = \frac{b}{|b|}$，此时 $\varphi = \frac{\pi}{2}$（当 $b>0$）或 $\frac{3\pi}{2}$（当 $b<0$）；若 $b=0$，则 $\sin\varphi=0$，$\cos\varphi = \frac{a}{|a|}$，此时 $\varphi = 0$（当 $a>0$）或 $\pi$（当 $a<0$）。
		\item[\textbf{\textcolor{CTrap}{错因分析（条件误解）：}}] 只记住公式形式，忽略条件“不同时为零”已包含一个为零的情形。
	\end{description}

	\medskip
	\begin{description}[style=nextline, leftmargin=2.8em, labelsep=0.5em, itemsep=0.35em]
		\item[\textbf{\textcolor{CTrap}{易错点二（象限判断错误）}}] 仅由 $\tan\varphi = \dfrac{b}{a}$ 求出 $\varphi$，不检验 $a,b$ 的符号，导致 $\varphi$ 取错象限。
		\item[\textbf{\textcolor{CTrap}{正确理解（看符号定象限）：}}] 必须同时看 $\cos\varphi = \dfrac{a}{\sqrt{a^{2}+b^{2}}}$ 和 $\sin\varphi = \dfrac{b}{\sqrt{a^{2}+b^{2}}}$ 的符号，从而确定 $\varphi$ 在哪个象限。
		\item[\textbf{\textcolor{CTrap}{错因分析（忽视正切周期）：}}] 正切函数周期为 $\pi$，仅由正切值无法区分 $\varphi$ 与 $\varphi+\pi$，必须依赖正余弦的正负。
	\end{description}

	\medskip
	\begin{description}[style=nextline, leftmargin=2.8em, labelsep=0.5em, itemsep=0.35em]
		\item[\textbf{\textcolor{CTrap}{易错点三（余弦形式混淆）}}] 使用余弦表达式 $a\sin x+b\cos x = \sqrt{a^{2}+b^{2}}\cos(x-\varphi')$ 时，误认为 $\varphi'$ 与 $\varphi$ 相同或关系不清。
		\item[\textbf{\textcolor{CTrap}{正确理解（区分两种形式）：}}] 余弦形式中 $\varphi'$ 满足 $\cos\varphi' = \dfrac{b}{\sqrt{a^{2}+b^{2}}}$，$\sin\varphi' = \dfrac{a}{\sqrt{a^{2}+b^{2}}}$，与正弦形式中的 $\varphi$ 相差 $\dfrac{\pi}{2}$，即 $\varphi' = \dfrac{\pi}{2} - \varphi$。
		\item[\textbf{\textcolor{CTrap}{错因分析（记混定义）：}}] 只记住正弦形式中的正余弦定义，直接套用到余弦形式导致符号交换错误。
	\end{description}
\end{trapbox}
